\uuid{TSCI}
\titre{Division euclidienne}
\auteur{bodin ; hanani ; mzari}
\organisation{exo7}
\licence{CC BY-NC-SA 4.0}
\datecreate{2020-01-20}
\niveau{L1}
\module{Algèbre}
\chapitre{Polynôme, fraction rationnelle}
\sousChapitre{Division euclidienne}
\difficulte{3}
\theme{Arithmétique des polynômes, division euclidienne}
\exo7qcmid{lille-1.polynomes.016}

\contenu{
\begin{question}
\qclassification{105.01,105.02}
\qtags{motcle=division euclidienne}
Soit $n\ge2$. Soit $A(X) = X^{2n}+X^{2n-2}$. Soit $B(X) = X^{n}+X^{n-1}$. Soit $A = BQ + R$ la division euclidienne de $A$ par $B$.
\begin{answers}
    \good{Le coefficient de $X^n$ de $Q$ est $1$.}

    \bad{Le coefficient de $X^{n-1}$ de $Q$ est $1$.}

    \good{Le coefficient de $X^{n-2}$ de $Q$ est $2$.}

    \good{$R$ est constitué d'un seul monôme.}

\end{answers}
\begin{explanations}
$Q(X) = X^n-X^{n-1}+2X^{n-2}-2X^{n-3}+\cdots$. $R(X) = \pm 2 X^{n-1}$.
\end{explanations}
\end{question}
}
